\documentclass{article}
\usepackage{graphicx}
\usepackage{listings}
\usepackage{color}

\definecolor{dkgreen}{rgb}{0,0.6,0}
\definecolor{gray}{rgb}{0.5,0.5,0.5}
\definecolor{mauve}{rgb}{0.58,0,0.82}

\lstset{frame=tb,
  language=Python,
  aboveskip=3mm,
  belowskip=3mm,
  showstringspaces=false,
  columns=flexible,
  basicstyle={\small\ttfamily},
  numbers=none,
  numberstyle=\tiny\color{gray},
  keywordstyle=\color{blue},
  commentstyle=\color{dkgreen},
  stringstyle=\color{mauve},
  breaklines=true,
  breakatwhitespace=true,
  tabsize=3
}

\title{Practical Work 1: TCP File Transfer}
\author{Vu Nhat Anh - 22BA13036}
\date{\today}

\begin{document}

\maketitle

\section{Protocol Design}
The protocol used for this file transfer is a simple TCP stream protocol.
\begin{itemize}
    \item \textbf{Transport Layer}: TCP is used to ensure reliable, ordered delivery of the data stream.
    \item \textbf{Connection}: The client initiates a connection to the server's IP and Port (8080).
    \item \textbf{Data Transfer}: 
    \begin{itemize}
        \item The client reads the file and sends raw bytes over the socket.
        \item The server receives the bytes and writes them directly to a file.
        \item The transfer ends when the client closes the connection (EOF detected by the server).
    \end{itemize}
\end{itemize}

\section{System Organization}
The system consists of two main components:
\begin{itemize}
    \item \textbf{Server}: Listens for incoming TCP connections. When a client connects, it accepts the connection, reads the data stream until the connection is closed, and saves the data to a file.
    \item \textbf{Client}: Connects to the server, reads a local file, and sends its content over the TCP connection.
\end{itemize}

\begin{figure}[h!]
    \centering
    % \includegraphics[width=0.8\textwidth]{system_diagram.png} 
    % Note: Diagram would be included here.
    \caption{Client-Server Architecture}
\end{figure}

\section{Implementation}

\subsection{Server Implementation}
The server uses the `socket` module to bind to a port and listen.
\begin{lstlisting}
with socket.socket(socket.AF_INET, socket.SOCK_STREAM) as s:
    s.bind((HOST, PORT))
    s.listen()
    conn, addr = s.accept()
    with conn:
        with open('received_file.txt', 'wb') as f:
            while True:
                data = conn.recv(1024)
                if not data: break
                f.write(data)
\end{lstlisting}

\subsection{Client Implementation}
The client connects to the server and sends the file data.
\begin{lstlisting}
with socket.socket(socket.AF_INET, socket.SOCK_STREAM) as s:
    s.connect((HOST, PORT))
    with open(filename, 'rb') as f:
        while True:
            data = f.read(1024)
            if not data: break
            s.sendall(data)
\end{lstlisting}

\section{Roles}
\begin{itemize}
    \item \textbf{Vu Nhat Anh}: Implemented both the Server and Client logic and wrote this report.
\end{itemize}

\end{document}
